\section{Die Nutzenarchitektur (INU/KNU/IDN)}
\label{sec:fepsde-utility}
% English: The Utility Architecture - INU/KNU/IDN Framework

%%%%%%%%%%%%%%%%%%%%%%%%%%%%%%%%%%%%%%%%%%%%%%%%%%%%%%%%%%%%%%%%%%%%%%%%%%%%%%%
% CHAPTER 10: Complete EBF Utility Framework
% Version 2.0 - Fully symmetric 144-component structure
%%%%%%%%%%%%%%%%%%%%%%%%%%%%%%%%%%%%%%%%%%%%%%%%%%%%%%%%%%%%%%%%%%%%%%%%%%%%%%%

\subsection{Die drei Nutzenkategorien: INU, KNU, IDN}
\label{subsec:three-categories}
% English: The Three Utility Categories

%%%%%%%%%%%%%%%%%%%%%%%%%%%%%%%%%%%%%%%%%%%%%%%%%%%%%%%%%%%%%%%%%%%%%%%%%%%%%%%
% 10.1: The theoretical foundation for the three utility categories
%%%%%%%%%%%%%%%%%%%%%%%%%%%%%%%%%%%%%%%%%%%%%%%%%%%%%%%%%%%%%%%%%%%%%%%%%%%%%%%

\subsubsection{The Problem with Single-Dimensional Utility}

Classical economics assumes a single utility function:
\begin{equation}
U_i = U_i(x_i)
\end{equation}
where person $i$ maximizes utility from their own consumption $x_i$.

This elegant formulation has powered economic theory for over a century.
But it cannot explain some of the most robust findings in behavioral economics:

\begin{itemize}
\item \textbf{Ultimatum Game:} People reject unfair offers, sacrificing money 
to punish unfairness \citep{PAP-guth1982experimental}
\item \textbf{Public Goods:} People contribute more than Nash equilibrium predicts, 
even without enforcement \citep{PAP-fehr2000cooperation}
\item \textbf{Third-Party Punishment:} People pay to punish unfairness 
even when they are not personally affected \citep{fehr2004third}
\item \textbf{Identity-Inconsistent Choices:} People reject higher-paying jobs 
that conflict with their self-image \citep{akerlofkranton}
\end{itemize}

These anomalies are not noise. They are systematic, replicable, and large.
They suggest that human motivation has \emph{multiple distinct sources}
that cannot be reduced to a single utility function.

\subsubsection{The Three Categories}

We propose that utility arises from three distinct categories,
each answering a different fundamental question:

\begin{definition}[The Three Utility Categories]
\label{def:three-categories}
\begin{align}
U^{INU} &: \text{\textbf{Individual Utility} --- ``What do I want for MYSELF?''} \\[0.5em]
U^{KNU} &: \text{\textbf{Collective Utility} --- ``What do I want for US?''} \\[0.5em]
U^{IDN} &: \text{\textbf{Identity Utility} --- ``Who AM I?''}
\end{align}
\end{definition}

These three categories are:
\begin{itemize}
\item \textbf{Conceptually distinct:} Each answers a different question
\item \textbf{Empirically separable:} Different experimental paradigms activate each
\item \textbf{Theoretically grounded:} Each connects to established research programs
\item \textbf{Structurally identical:} All three have the same internal architecture (Section~\ref{subsec:144-structure})
\end{itemize}

%%%%%%%%%%%%%%%%%%%%%%%%%%%%%%%%%%%%%%%%%%%%%%%%%%%%%%%%%%%%%%%%%%%%%%%%%%%%%%%
\subsubsection{Individual Utility (INU)}
%%%%%%%%%%%%%%%%%%%%%%%%%%%%%%%%%%%%%%%%%%%%%%%%%%%%%%%%%%%%%%%%%%%%%%%%%%%%%%%

\paragraph{Definition.}
Individual Utility captures preferences over personal outcomes:
my income, my health, my relationships, my experiences.

\begin{equation}
U^{INU} = U^{INU}(x_{self})
\end{equation}

This is the utility that Arrow-Debreu models, that revealed preference measures,
that cost-benefit analyses aggregate. It is real, important, and well-understood.

\paragraph{Theoretical Foundation.}
\begin{itemize}
\item Classical utility theory \citep{vonneumann1944}
\item Revealed preference \citep{samuelson1938}
\item Expected utility \citep{savage1954}
\item Prospect theory (for reference-dependent INU) \citep{kahnemantversky1979}
\end{itemize}

\paragraph{The Question INU Answers.}
\begin{quote}
``What consequences does this action have \emph{for me personally}?''
\end{quote}

\paragraph{When INU Dominates.}
INU tends to dominate decision-making when:
\begin{itemize}
\item Consequences are primarily personal
\item No salient group membership is activated
\item Identity is not threatened
\item Time pressure forces heuristic processing (System 1)
\end{itemize}

\paragraph{What INU Cannot Explain.}
INU alone cannot explain:
\begin{itemize}
\item Rejecting \euro{1} in ultimatum game (INU says: \euro{1} $>$ \euro{0})
\item Voting (INU says: costs exceed expected benefits)
\item Anonymous charitable giving (INU says: no personal return)
\item Refusing lucrative jobs that ``aren't me'' (INU says: more money = better)
\end{itemize}

%%%%%%%%%%%%%%%%%%%%%%%%%%%%%%%%%%%%%%%%%%%%%%%%%%%%%%%%%%%%%%%%%%%%%%%%%%%%%%%
\subsubsection{Collective Utility (KNU)}
%%%%%%%%%%%%%%%%%%%%%%%%%%%%%%%%%%%%%%%%%%%%%%%%%%%%%%%%%%%%%%%%%%%%%%%%%%%%%%%

\paragraph{Definition.}
Collective Utility captures preferences over group outcomes:
my family's welfare, my community's prosperity, my nation's success,
my company's performance---weighted by fairness and distribution.

\begin{equation}
U^{KNU} = U^{KNU}(x_{group}, \text{fairness}, \text{distribution})
\end{equation}

\paragraph{Theoretical Foundation.}
The existence of KNU as a distinct utility category is established by
decades of experimental research, much of it from the Zurich tradition:

\begin{itemize}
\item \textbf{Inequity Aversion:} \citet{PAP-fehr1999theory} showed that people 
dislike both advantageous and disadvantageous inequality, and will sacrifice 
personal payoffs to reduce it.

\item \textbf{Reciprocity:} \citet{PAP-falk2006theory} demonstrated that people 
reward kind intentions and punish unkind ones, even at personal cost.

\item \textbf{Social Preferences:} \citet{charness2002understanding} found that 
people care about both efficiency (total welfare) and distribution (fairness).

\item \textbf{Conditional Cooperation:} \citet{fischbacher2001people} showed that 
most people are ``conditional cooperators'' who contribute to public goods 
if others do too.
\end{itemize}

\paragraph{The Question KNU Answers.}
\begin{quote}
``What consequences does this action have \emph{for my group}?
Is the outcome \emph{fair}?''
\end{quote}

\paragraph{Why KNU Is Distinct from INU.}
KNU is not reducible to ``enlightened self-interest'' or ``warm glow.''
The evidence shows:

\begin{enumerate}
\item \textbf{People sacrifice for groups they won't interact with again.}
One-shot anonymous games still show cooperation and punishment.

\item \textbf{The magnitude is too large for reputation effects.}
Third-party punishment costs can exceed 20\% of endowment.

\item \textbf{Fairness matters independently of outcomes.}
People reject unfair processes even when outcomes would be identical.

\item \textbf{Group identity modulates the effect.}
KNU is stronger for in-groups than out-groups \citep{chen2009group}.
\end{enumerate}

\paragraph{When KNU Dominates.}
KNU tends to dominate when:
\begin{itemize}
\item Group membership is salient (``we'' language, group primes)
\item Fairness norms are activated
\item Others' behavior is visible (conditional cooperation)
\item Social observation increases accountability
\end{itemize}

\paragraph{KNU Is Level-Dependent.}
A critical insight: ``we'' depends on \emph{which group} is salient.
KNU is not a single value but a \textbf{level-indexed function}:

\begin{equation}
U^{KNU,L} = U^{KNU}(x_{group}^L, \text{fairness}^L, \text{distribution}^L)
\label{eq:knu-level}
\end{equation}

where $L \in \{1, 2, 3, 4, 5, 6, \infty\}$ corresponds to the aggregation hierarchy
from Appendix~AAA (CORE-WHO):

\begin{center}
\renewcommand{\arraystretch}{1.2}
\begin{tabular}{clll}
\toprule
\textbf{Level} & \textbf{``We''} & \textbf{KNU Concern} & \textbf{Example} \\
\midrule
$L = 1$ & Dyad & Partner's welfare & ``Is this fair to my spouse?'' \\
$L = 2$ & Family & Household welfare & ``Can we afford this?'' \\
$L = 3$ & Tribe & Community welfare & ``What will the neighbors think?'' \\
$L = 4$ & Organization & Firm/team welfare & ``Is this good for the company?'' \\
$L = 5$ & Nation & National welfare & ``Is this good for the country?'' \\
$L = 6$ & Civilization & Human welfare & ``Is this good for humanity?'' \\
$L = \infty$ & Meta & All beings & ``Is this good for the planet?'' \\
\bottomrule
\end{tabular}
\end{center}

\paragraph{Level Salience.}
Which level dominates depends on context via the \textbf{level salience function}
$\alpha^L(\Psi)$ (see Appendix~AAA, Section~3):

\begin{equation}
KNU_{total} = \sum_{L=1}^{\infty} \alpha^L(\Psi) \cdot U^{KNU,L}
\label{eq:knu-aggregation}
\end{equation}

where $\sum_L \alpha^L(\Psi) = 1$ and $\alpha^L(\Psi)$ is determined by:
\begin{itemize}[nosep]
\item $\Psi_S$ (Social context): High trust $\rightarrow$ higher-level KNU salient
\item $\Psi_I$ (Institutional context): Strong institutions $\rightarrow$ formal levels (4, 5) salient
\item $\Psi_C$ (Cognitive load): High load $\rightarrow$ proximate levels (1, 2, 3) dominate
\end{itemize}

\textit{Cross-Reference:} The complete level hierarchy and salience functions are
formalized in Appendix~AAA (CORE-WHO: Who has utility?).

\paragraph{Implication for the 9C Framework.}
This level-indexing reveals a deep integration within the 9C CORE architecture:
\textbf{WHO} (which level $L$) determines how collective utility is structured,
\textbf{WHEN} (context $\Psi$) determines which level is salient via $\alpha^L(\Psi)$,
and together they feed into \textbf{WHAT} (the utility dimensions).

\begin{tcolorbox}[colback=yellow!5!white, colframe=yellow!75!black,
    title=\textbf{The Same Decision, Different Level Activation}]
Consider accepting a job offer in another city:

\begin{center}
\renewcommand{\arraystretch}{1.2}
\begin{tabular}{lll}
\toprule
\textbf{Context} & \textbf{Salient Level} & \textbf{KNU Question} \\
\midrule
Family dinner table & $L=2$ (Family) & ``Can WE afford this move?'' \\
Team meeting at work & $L=4$ (Organization) & ``Will this hurt the TEAM?'' \\
National policy debate & $L=5$ (Nation) & ``Is this brain drain for the COUNTRY?'' \\
Climate discussion & $L=\infty$ (Meta) & ``What's the carbon footprint?'' \\
\bottomrule
\end{tabular}
\end{center}

The \emph{objective} decision is identical. But the \emph{KNU evaluation} differs
because context activates different levels. This is why behavioral predictions
require knowing not just WHAT people value, but WHO they identify with in the moment.
\end{tcolorbox}

\paragraph{KNU Examples.}

\begin{center}
\renewcommand{\arraystretch}{1.3}
\begin{tabular}{p{4cm}p{4cm}p{4cm}}
\toprule
\textbf{Situation} & \textbf{INU Prediction} & \textbf{KNU Prediction} \\
\midrule
Ultimatum: offered \euro{1}/\euro{10} & Accept (\euro{1} $>$ \euro{0}) & Reject (unfair!) \\
Public good: contribute? & Free-ride & Contribute if others do \\
Bonus: self vs. team? & Self & Team (if identified) \\
Tax compliance & Evade if undetected & Pay (fair share) \\
\bottomrule
\end{tabular}
\end{center}

%%%%%%%%%%%%%%%%%%%%%%%%%%%%%%%%%%%%%%%%%%%%%%%%%%%%%%%%%%%%%%%%%%%%%%%%%%%%%%%
\subsubsection{Identity Utility (IDN)}
%%%%%%%%%%%%%%%%%%%%%%%%%%%%%%%%%%%%%%%%%%%%%%%%%%%%%%%%%%%%%%%%%%%%%%%%%%%%%%%

\paragraph{Definition.}
Identity Utility captures preferences over self-concept:
being a certain kind of person, living according to one's values,
maintaining a coherent life narrative, acting consistently with who one is.

\begin{equation}
U^{IDN} = U^{IDN}(\text{match between action and self-concept})
\end{equation}

\paragraph{Theoretical Foundation.}

\begin{itemize}
\item \textbf{Identity Economics:} \citet{akerlofkranton} formalized how 
identity enters utility. People have identity $I$, which prescribes ideal 
behavior $a^*$. Deviations from $a^*$ create disutility.

\item \textbf{Self-Signaling:} \citet{benaboutirole} showed that people 
take actions partly to learn about (or maintain beliefs about) their own type.
``I donate to charity, therefore I am generous.''

\item \textbf{Moral Identity:} \citet{aquino2002self} demonstrated that 
people for whom morality is central to identity behave more ethically.

\item \textbf{Professional Identity:} Extensive literature shows that 
occupational identity shapes behavior independently of pay \citep{ashforth1989social}.
\end{itemize}

\paragraph{The Question IDN Answers.}
\begin{quote}
``Is this action consistent with \emph{who I am}?
What does this choice \emph{say about me}?''
\end{quote}

\paragraph{Why IDN Is Distinct from INU and KNU.}
People sacrifice both personal benefit (INU) and group benefit (KNU) for identity:

\begin{enumerate}
\item \textbf{Refusing to lie even when profitable and harmless.}
Neither INU (profit) nor KNU (no one harmed) can explain this.
But IDN can: ``I am not a liar.''

\item \textbf{Maintaining principles against group pressure.}
When the group wants you to compromise values, KNU says comply.
IDN says resist: ``I have to live with myself.''

\item \textbf{Choosing meaningful work over higher pay.}
INU says take the money. IDN says: ``That's not who I am.''

\item \textbf{Coal miners refusing to become coders.}
INU says: higher pay, better conditions, take it.
IDN says: ``I'm a miner. That's who I am. Who would I be if I sat at a computer?''
\end{enumerate}

\paragraph{The Identity Utility Function.}
Following \citet{akerlofkranton}, identity utility depends on the 
\emph{match} between actual state and identity-prescribed state:

\begin{equation}
U^{IDN}_d = -t_d \cdot |x_d - x^*_d|
\label{eq:identity-match}
\end{equation}

where:
\begin{itemize}
\item $x_d$ = actual state in dimension $d$
\item $x^*_d$ = identity-prescribed state (``someone like me should have/be...'')
\item $t_d$ = identity salience for dimension $d$
\item The negative sign indicates disutility from mismatch
\end{itemize}

\paragraph{When IDN Dominates.}
IDN tends to dominate when:
\begin{itemize}
\item Identity is threatened or challenged
\item Choice has symbolic meaning (``what does this say about me?'')
\item Public visibility increases self-presentation concerns
\item Transitions threaten core roles (career change, retirement, divorce)
\end{itemize}

\paragraph{IDN Examples.}

\begin{center}
\renewcommand{\arraystretch}{1.3}
\begin{tabular}{p{4cm}p{3.5cm}p{4.5cm}}
\toprule
\textbf{Situation} & \textbf{INU/KNU Prediction} & \textbf{IDN Prediction} \\
\midrule
Coal miner $\rightarrow$ coder (+\$20k) & Accept (more money) & Reject (not who I am) \\
Vegetarian offered steak & Accept if tasty & Reject (I don't eat meat) \\
Lie for small profit, no victim & Tell lie (profit, no harm) & Don't lie (I'm honest) \\
Academic takes corporate job & Accept if pays more & Reject (I'm a scholar) \\
\bottomrule
\end{tabular}
\end{center}

%%%%%%%%%%%%%%%%%%%%%%%%%%%%%%%%%%%%%%%%%%%%%%%%%%%%%%%%%%%%%%%%%%%%%%%%%%%%%%%
\subsubsection{The Relationship Between the Three Categories}
%%%%%%%%%%%%%%%%%%%%%%%%%%%%%%%%%%%%%%%%%%%%%%%%%%%%%%%%%%%%%%%%%%%%%%%%%%%%%%%

\paragraph{Not Additive, but Interactive.}
The three utility categories are not simply added.
They interact through complementarities:

\begin{equation}
Q = \sum_{i \in \{INU, KNU, IDN\}} \omega_i \cdot U^i + \sum_{i \neq j} \gamma_{ij} \cdot U^i \cdot U^j
\label{eq:utility-complementarity}
\end{equation}

where $\gamma_{ij}$ captures the interaction between categories.

\paragraph{Complementarity Patterns.}

\begin{table}[htbp]
\centering
\caption{Inter-Category Complementarities}
\label{tab:complementarities}
\renewcommand{\arraystretch}{1.3}
\begin{tabular}{@{}lcll@{}}
\toprule
\textbf{Interaction} & \textbf{Sign} & \textbf{Meaning} & \textbf{Example} \\
\midrule
$\gamma_{INU,KNU}$ & + & My gain = group gain & Team bonus, cooperative \\
$\gamma_{INU,KNU}$ & $-$ & My gain = group loss & Corruption, free-riding \\
\addlinespace
$\gamma_{INU,IDN}$ & + & Success confirms identity & ``I earned this'' \\
$\gamma_{INU,IDN}$ & $-$ & Success violates identity & ``Selling out'' \\
\addlinespace
$\gamma_{KNU,IDN}$ & + & Group identity & ``We are miners'' \\
$\gamma_{KNU,IDN}$ & $-$ & Group pressure vs. self & Conformity vs. authenticity \\
\bottomrule
\end{tabular}
\end{table}

\paragraph{The Coherence Condition.}
Maximum welfare occurs when all three categories align:
\begin{equation}
K^{utility} = \text{Alignment}(U^{INU}, U^{KNU}, U^{IDN})
\end{equation}

\begin{table}[htbp]
\centering
\caption{Utility Coherence States}
\label{tab:utility-coherence}
\renewcommand{\arraystretch}{1.3}
\begin{tabular}{@{}cccll@{}}
\toprule
$U^{INU}$ & $U^{KNU}$ & $U^{IDN}$ & \textbf{Coherence} & \textbf{Example} \\
\midrule
+ & + & + & HIGH & Successful entrepreneur helping community \\
+ & $-$ & ? & LOW & Profitable but exploitative \\
+ & + & $-$ & LOW & Everyone wins but ``this isn't me'' \\
$-$ & $-$ & + & MEDIUM & Poor but proud, ``we suffer together'' \\
\bottomrule
\end{tabular}
\end{table}

%%%%%%%%%%%%%%%%%%%%%%%%%%%%%%%%%%%%%%%%%%%%%%%%%%%%%%%%%%%%%%%%%%%%%%%%%%%%%%%
\subsubsection{Structural Identity: All Three Have the Same Architecture}
%%%%%%%%%%%%%%%%%%%%%%%%%%%%%%%%%%%%%%%%%%%%%%%%%%%%%%%%%%%%%%%%%%%%%%%%%%%%%%%

\paragraph{The Key Insight.}
While INU, KNU, and IDN answer different questions, they share 
\emph{identical internal structure}:

\begin{tcolorbox}[colback=blue!5!white, colframe=blue!75!black, title=Structural Identity]
Each category (INU, KNU, IDN) has:
\begin{itemize}
\item The same 6 dimensions (FEPSDE)
\item The same 4 time horizons
\item The same Gain/Pain structure
\item Its own reference points
\item Its own $\lambda$ (loss aversion) parameters
\item Its own $\delta$ (discount) parameters
\end{itemize}

The difference is not in \emph{structure} but in \emph{parameters}---specifically, 
how reference points are formed.
\end{tcolorbox}

This structural identity is what enables the clean 144-component architecture
developed in Section~\ref{subsec:144-structure}.

\paragraph{Reference Point Formation: The Key Difference.}

\begin{table}[htbp]
\centering
\caption{Reference Point Formation by Category}
\label{tab:reference-points}
\renewcommand{\arraystretch}{1.4}
\begin{tabular}{@{}llp{6cm}@{}}
\toprule
\textbf{Category} & \textbf{Reference Point} & \textbf{Formation Process} \\
\midrule
INU & $r^{INU}_d$ & Status quo, expectations, social comparison \\
& & ``What do I have? What did I expect? What do peers have?'' \\
\addlinespace
KNU & $r^{KNU}_d$ & Group norm, fair distribution, collective standard \\
& & ``What should we have? What is fair? What do similar groups have?'' \\
\addlinespace
IDN & $r^{IDN}_d$ & Self-concept, identity standard, role expectation \\
& & ``What should someone like me have/be? What does my identity prescribe?'' \\
\bottomrule
\end{tabular}
\end{table}

\paragraph{Example: Financial Dimension.}
Consider a manager earning \euro{100,000} who loses their job and finds 
a new one at \euro{60,000}.

\begin{table}[htbp]
\centering
\caption{Reference Points in Action: Job Loss Example}
\label{tab:reference-example}
\renewcommand{\arraystretch}{1.3}
\begin{tabular}{@{}llccc@{}}
\toprule
\textbf{Category} & \textbf{Reference Point} & \textbf{Value} & \textbf{Outcome} & \textbf{Pain} \\
\midrule
INU & Status quo (old salary) & \euro{100k} & \euro{60k} & \euro{40k} \\
KNU & Family standard & \euro{80k} & \euro{60k} & \euro{20k} \\
IDN & ``Successful manager'' & \euro{120k} & \euro{60k} & \euro{60k} \\
\bottomrule
\end{tabular}
\end{table}

The \emph{same} objective outcome (\euro{60k} salary) creates different 
pains depending on the reference point, which differs by category.

%%%%%%%%%%%%%%%%%%%%%%%%%%%%%%%%%%%%%%%%%%%%%%%%%%%%%%%%%%%%%%%%%%%%%%%%%%%%%%%
\subsubsection{Summary: The Three Utility Categories}
%%%%%%%%%%%%%%%%%%%%%%%%%%%%%%%%%%%%%%%%%%%%%%%%%%%%%%%%%%%%%%%%%%%%%%%%%%%%%%%

\begin{tcolorbox}[colback=gray!5!white, colframe=gray!75!black, title=Section 10.1 Summary]
\textbf{Three distinct utility categories:}
\begin{itemize}
\item \textbf{INU (Individual):} What do I want for myself?
\item \textbf{KNU (Collective):} What do I want for us?
\item \textbf{IDN (Identity):} Who am I?
\end{itemize}

\textbf{Key properties:}
\begin{itemize}
\item Conceptually distinct (different questions)
\item Empirically separable (different experiments activate each)
\item Theoretically grounded (classical econ, behavioral econ, identity econ)
\item \textbf{Structurally identical} (same FEPSDE $\times$ Time $\times$ Gain/Pain architecture)
\item Different only in reference point formation
\end{itemize}

\textbf{Interaction:}
\begin{itemize}
\item Categories interact through complementarities ($\gamma_{ij}$)
\item Maximum welfare when all three align
\item Conflicts create psychological tension and behavioral instability
\end{itemize}

\textbf{Foundation for:}
\begin{itemize}
\item 144-component structure (Section~\ref{subsec:144-structure})
\item Awareness function (Chapter 11)
\item Willingness to change (Chapter 12)
\item Behavioral change segments (Chapter 14)
\end{itemize}
\end{tcolorbox}

\vspace{1em}
\hrule
\vspace{0.5em}
\noindent\textit{Next section: 10.2 The Six FEPSDE Dimensions}

%%%%%%%%%%%%%%%%%%%%%%%%%%%%%%%%%%%%%%%%%%%%%%%%%%%%%%%%%%%%%%%%%%%%%%%%%%%%%%%
